\documentclass[12pt, a4paper]{article}

%=============== PAQUETES ===============
\usepackage[utf8]{inputenc}
\usepackage[T1]{fontenc}
\usepackage[spanish,es-tabla]{babel} % es-tabla uses "Tabla" instead of "Cuadro"
\usepackage{geometry}
\usepackage{amsmath}
\usepackage{graphicx}
\usepackage{xcolor}
\usepackage{circuitikz}
\usetikzlibrary{arrows.meta, babel}
\usepackage{comment}

% Paquetes adicionales para mejor formato
\usepackage{siunitx}
\usepackage{enumitem}
\usepackage{float}
\usepackage{booktabs}
\usepackage{fancyhdr}
\usepackage{hyperref} % Cargar al final (generalmente)

%=============== CONFIGURACIÓN ===============
\geometry{a4paper, margin=2.5cm}

% Configuración de unidades SI
\sisetup{
    output-decimal-marker = {.}, % Usar punto decimal como en el original
    range-phrase = { a },
    range-units = single,
    per-mode = symbol
}

% Configuración de enlaces
\hypersetup{
    colorlinks=true,
    linkcolor=blue,
    urlcolor=cyan,
    pdftitle={Práctica Semana 7: Diodo Zener, Transitorios y Pequeña Señal},
}

% Configuración de encabezado y pie de página
\pagestyle{fancy}
\fancyhf{}
\setlength{\headheight}{14.5pt} % Fix for fancyhdr warning
\lhead{Electrónica - Dispositivos}
\rhead{Práctica Semana 7}
\cfoot{\thepage}

\newcommand{\code}[1]{\texttt{#1}}

%=============== DOCUMENTO ===============

\begin{document}

\title{Práctica de Laboratorio Semana 7 \\ \large El Diodo: Zener, Transitorios y Pequeña Señal}
\author{Plan de Estudios de Electrónica}
\date{\today}
\maketitle

\section*{Objetivos}

\begin{itemize}[noitemsep]
    \item Analizar y diseñar un circuito regulador de tensión básico utilizando un diodo Zener.
    \item Comprender el comportamiento transitorio de un diodo, observando los efectos de la capacidad de la juntura.
    \item Aplicar el modelo de pequeña señal para determinar la resistencia dinámica de un diodo en un punto de operación DC.
    \item Utilizar un simulador de circuitos (SPICE, LTspice, DroidTesla, etc.) para verificar los cálculos teóricos y visualizar el comportamiento de los circuitos.
\end{itemize}

\section*{Introducción Teórica}

Esta práctica se centra en tres aspectos avanzados del diodo.
\begin{enumerate}
    \item \textbf{Diodo Zener:} Es un diodo diseñado para operar en la región de ruptura inversa de forma controlada. Mantiene una tensión casi constante a través de sus terminales ($V_Z$) para un rango de corrientes inversas, lo que lo hace ideal para aplicaciones de regulación de voltaje.
    \item \textbf{Comportamiento Transitorio:} Un diodo no conmuta instantáneamente. La \textbf{capacidad de difusión} (dominante en polarización directa) y la \textbf{capacidad de transición o de juntura} (dominante en polarización inversa) introducen retardos. Estos efectos son cruciales en aplicaciones de alta frecuencia.
    \item \textbf{Modelo de Pequeña Señal:} Cuando un diodo está polarizado en un punto de operación DC (Q-point) y se le superpone una pequeña señal AC, su comportamiento para esa señal AC puede ser modelado por una resistencia lineal llamada \textbf{resistencia dinámica} ($r_d$), cuyo valor es $r_d = V_T / I_{DQ}$, donde $V_T$ es la tensión térmica (aprox. \SI{25}{\milli\volt} a temperatura ambiente) e $I_{DQ}$ es la corriente DC que atraviesa el diodo.
\end{enumerate}

\section*{Materiales y Equipo}

\begin{itemize}[noitemsep]
    \item Software de simulación de circuitos (ej. DroidTesla, LTspice, Multisim, o cualquier variante de SPICE).
    \item Componentes para simulación:
    \begin{itemize}[noitemsep]
        \item Diodo Zener (ej. 1N4733A, $V_Z=\SI{5.1}{\volt}$).
        \item Diodo de señal (ej. 1N4148).
        \item Resistencias de varios valores.
        \item Fuentes de tensión DC y AC (sinusoidal, pulso).
    \end{itemize}
\end{itemize}

\newpage

\section*{Procedimiento Experimental}

\subsection*{Parte 1: Regulador de Tensión con Diodo Zener}

Considere el siguiente circuito regulador. El objetivo es mantener una tensión estable de \SI{5.1}{\volt} en la carga $R_L$.

\begin{center}
\begin{circuitikz}[scale=1.2, transform shape] % Circuito mejorado para mayor claridad
    \draw (0,2) to[sV, l=$V_{in}$] (0,0) node[ground]{};
    \draw (0,2) to[R, l=$R_S$] (3,2);
    \draw (3,2) to[zD, *-*, l_=$D_Z$] (3,0) node[ground]{};
    \draw (3,2) to[R, l=$R_L$] (5,2) to[short] (5,0) node[ground]{};
\end{circuitikz}
\end{center}

\begin{enumerate}
    \item \textbf{Diseño:} Use un diodo Zener 1N4733A ($V_Z \approx \SI{5.1}{\volt}$). Si la tensión de entrada $V_{in}$ varía entre \SI{10}{\volt} y \SI{15}{\volt} y la carga $R_L$ es de \SI{1}{\kilo\ohm}, calcule un valor adecuado para $R_S$ que garantice que el Zener siempre esté en su región de regulación. Asegúrese de que la potencia disipada por el Zener no exceda su límite.

    \item \textbf{Simulación (Regulación de Línea):}
    \begin{itemize}
        \item Monte el circuito en el simulador con los valores calculados.
        \item Realice un análisis DC (\code{.dc}) variando $V_{in}$ de \SI{0}{\volt} a \SI{20}{\volt}.
        \item Grafique la tensión de salida $V_{out}$ (en $R_L$) en función de $V_{in}$.
        \item Observe el punto donde $V_{out}$ se estabiliza en $V_Z$. ¿Coincide con sus cálculos?
    \end{itemize}

    \item \textbf{Simulación (Regulación de Carga):}
    \begin{itemize}
        \item Fije $V_{in}$ a un valor constante (ej. \SI{12}{\volt}).
        \item Realice un análisis DC variando el valor de $R_L$ (ej. desde \SI{100}{\ohm} hasta \SI{5}{\kilo\ohm}).
        \item Grafique $V_{out}$ en función de $R_L$.
        \item Identifique el valor mínimo de $R_L$ para el cual el circuito sigue regulando.
    \end{itemize}
\end{enumerate}

\subsection*{Parte 2: Análisis Transitorio del Diodo}

Se analizará el tiempo de recuperación inversa de un diodo 1N4148.

\begin{center}
\begin{circuitikz}[scale=1.2, transform shape]
    % Fuente AC
    \draw (0,0) node[ground]{} 
        to[sV, l=$V_{in}$] (0,2) 
        
    % Resistencia (con formato de texto correcto para la unidad)
    % Usamos \text{k} para que no salga cursiva
        to[R, l={$R_1=\SI{1}{\kilo\ohm}$}] (3,2) 
        
    % Diodo (etiqueta en modo texto normal es correcta)
        to[D, l=1N4148] (3,0) 
        node[ground]{};
        
    % Opcional: Si quieres visualizar la salida Vout sobre el diodo
    \draw (3,2) to[short, -o] (4,2) node[right] {$V_{out}$};
\end{circuitikz}
\end{center}

\begin{enumerate}
    \item \textbf{Configuración:} Configure la fuente $V_{in}$ como una señal de pulso (\code{PULSE}) que alterne entre \SI{-5}{\volt} y \SI{5}{\volt}, con un período de \SI{1}{\micro\second} y tiempos de subida/bajada muy cortos (ej. \SI{1}{\nano\second}).
    
    \item \textbf{Simulación:} Realice un análisis transitorio (\code{.tran}) durante unos pocos ciclos (ej. \SI{3}{\micro\second}).
    
    \item \textbf{Análisis de Resultados:}
    \begin{itemize}
        \item Grafique la tensión en el diodo ($V_D$) y la corriente a través de él ($I_D$).
        \item Observe el comportamiento cuando la entrada cambia de +5V a -5V. Notará que el diodo no se apaga instantáneamente y conduce una corriente inversa por un corto período. Este es el \textbf{tiempo de almacenamiento} debido a la carga almacenada en la juntura.
        \item Mida este tiempo de recuperación inversa en la gráfica.
    \end{itemize}
\end{enumerate}

\subsection*{Parte 3: Modelo de Pequeña Señal}

Se verificará el concepto de resistencia dinámica ($r_d$).

\begin{center}
\begin{circuitikz}[scale=1.2, transform shape]
    % Rama Izquierda: Fuentes en Serie
    % Dividí la altura total (3) en dos partes iguales (1.5 cada una)
    \draw (0,0) node[ground]{} 
        to[V, l={$V_{DC}=\SI{10}{\volt}$}] (0,1.5)  % Fuente DC (Batería)
        to[sV, l=$v_{ac}$] (0,3)                % Fuente AC (Sinusoidal)
        
    % Rama Superior: Resistencia
    % Usamos \text{k} para la unidad correcta
        to[R, l={$R_1=\SI{1}{\kilo\ohm}$}] (3,3) 
        
    % Rama Derecha: Diodo
        to[D, l=1N4148] (3,0) 
        node[ground]{};

    % (Opcional) Terminal de salida para mayor claridad
    \draw (3,3) to[short, -o] (4,3) node[right] {$V_{out}$};
\end{circuitikz}
\end{center}

\begin{enumerate}
    \item \textbf{Cálculo Teórico:}
    \begin{itemize}
        \item Primero, ignore la fuente AC ($v_{ac}$). Calcule la corriente DC ($I_{DQ}$) que fluye a través del diodo, asumiendo una caída de tensión de \SI{0.7}{\volt} en él.
        \item Con $I_{DQ}$, calcule la resistencia dinámica teórica: $r_d = V_T / I_{DQ}$ (use $V_T \approx \SI{25}{\milli\volt}$).
    \end{itemize}

    \item \textbf{Simulación:}
    \begin{itemize}
        \item Monte el circuito. Configure $v_{ac}$ como una fuente sinusoidal (\code{SINE}) con \SI{0}{\volt} de offset, \SI{1}{\milli\volt} de amplitud y \SI{1}{\kilo\hertz} de frecuencia.
        \item Realice un análisis transitorio (\code{.tran}) para visualizar las señales.
        \item Grafique la corriente en el diodo ($I_D$) y la tensión en el diodo ($V_D$). Verá una pequeña ondulación AC superpuesta sobre los valores DC.
        \item Mida la amplitud pico a pico de la corriente AC ($\Delta i_d$) y de la tensión AC ($\Delta v_d$) en el diodo.
        \item Calcule la resistencia dinámica a partir de la simulación: $r_{d,sim} = \Delta v_d / \Delta i_d$.
    \end{itemize}

    \item \textbf{Comparación:} Compare el valor de $r_d$ teórico con el obtenido en la simulación ($r_{d,sim}$). Discuta cualquier discrepancia.
\end{enumerate}

\section*{Preguntas de Cierre}
\begin{enumerate}
    \item En la Parte 1, ¿qué sucede si la carga $R_L$ es demasiado pequeña? ¿Por qué el regulador deja de funcionar?
    \item En la Parte 2, ¿por qué es importante el tiempo de recuperación inversa en el diseño de fuentes de alimentación conmutadas (SMPS)?
    \item En la Parte 3, ¿qué pasaría con el valor de $r_d$ si se aumentara la tensión de la fuente $V_{DC}$ a \SI{20}{\volt}? Justifique su respuesta.
\end{enumerate}

\end{document}